\chapter{Reference: tabs}
\label{ch:refman-tabs}

\begin{aside}
Tabs partition a single content area into named panels.
Store the active index in the Model; render the matching
content; show the labels above.
\end{aside}

\section{Working example}

\begin{lstlisting}[language=C++]
#include <maya/maya.hpp>
#include <vector>

using namespace maya;
using namespace maya::dsl;

struct TabsDemo {
    struct Model {
        std::vector<std::string> labels{"info", "logs", "about"};
        int active = 0;
    };
    struct Next {};
    struct Prev {};
    struct Pick { int i; };
    struct Quit {};
    using Msg = std::variant<Next, Prev, Pick, Quit>;

    static Model init() { return {}; }

    static auto update(Model m, Msg msg)
        -> std::pair<Model, Cmd<Msg>>
    {
        int n = static_cast<int>(m.labels.size());
        return std::visit(overload{
            [&](Next) { m.active = (m.active + 1) % n; return std::pair{m, Cmd<Msg>{}}; },
            [&](Prev) { m.active = (m.active - 1 + n) % n; return std::pair{m, Cmd<Msg>{}}; },
            [&](Pick p) {
                if (p.i >= 0 && p.i < n) m.active = p.i;
                return std::pair{m, Cmd<Msg>{}};
            },
            [](Quit) {
                return std::pair{Model{}, Cmd<Msg>::quit()};
            },
        }, msg);
    }

    static Element view(const Model& m) {
        std::vector<Element> tabs;
        for (size_t i = 0; i < m.labels.size(); ++i) {
            auto label = text(" " + m.labels[i] + " ");
            if (static_cast<int>(i) == m.active) {
                tabs.push_back(label | Bold | Bg<60, 100, 180>);
            } else {
                tabs.push_back(label | Dim);
            }
        }

        Element body;
        switch (m.active) {
            case 0: body = t<"info: hello, world.">; break;
            case 1: body = t<"logs: nothing to report.">; break;
            default: body = t<"about: built with maya.">; break;
        }

        return v(
            h(tabs),
            blank_,
            body,
            blank_,
            t<"tab/shift-tab cycle, 1-3 pick, q quit"> | Dim
        ) | pad<1> | border_<Round>;
    }

    static auto subscribe(const Model&) -> Sub<Msg> {
        return key_map<Msg>({
            {'q',                Quit{}},
            {SpecialKey::Tab,    Next{}},
            {shift(SpecialKey::Tab), Prev{}},
            {'1', Pick{0}},
            {'2', Pick{1}},
            {'3', Pick{2}},
        });
    }
};

int main() { run<TabsDemo>({.title = "tabs"}); }
\end{lstlisting}

\section{Notes}

\begin{itemize}[leftmargin=*]
  \item Use a clear visual distinction between active and
    inactive tabs (background colour beats underline).
  \item Always include keyboard shortcuts; mouse-only tabs
    fail half your audience.
  \item For many tabs, scroll the label row horizontally.
\end{itemize}

\section{Recap}

\begin{recap}
\begin{itemize}[leftmargin=*]
  \item One active index; one switch in render.
  \item Tab to cycle; digits to jump.
  \item Bold + background highlights the active tab.
\end{itemize}
\end{recap}

\section*{Workshop}

\begin{enumerate}[leftmargin=*]
  \item Wrap each tab body in a separate function.
  \item Add a close-tab key for dynamic tab sets.
  \item Persist the active tab across runs.
\end{enumerate}
